\documentclass[11pt]{article}

\usepackage[T1]{fontenc}
\usepackage{amsmath,amssymb,amsthm}
\usepackage{booktabs}
\usepackage{geometry}
\usepackage{hyperref}
\usepackage{xurl}

\geometry{margin=1in}
\hypersetup{
  colorlinks=true,
  linkcolor=blue,
  citecolor=blue,
  urlcolor=blue,
  pdftitle={A Certified Order-Twelve Extension of the gamma--theta Frontier in One-Guard Eternal Domination},
  pdfauthor={Alec Kriebel},
  pdfsubject={One-guard eternal domination and clique covering},
  pdfkeywords={eternal domination, domination number, clique cover, gamma--theta conjecture}
}

\newtheorem{theorem}{Theorem}[section]
\newtheorem{lemma}[theorem]{Lemma}
\newtheorem{proposition}[theorem]{Proposition}
\newtheorem{corollary}[theorem]{Corollary}
\theoremstyle{definition}
\newtheorem{definition}[theorem]{Definition}
\theoremstyle{remark}
\newtheorem{remark}[theorem]{Remark}

\newcommand{\gaminf}{\gamma^{\infty}}
\newcommand{\V}{V}
\newcommand{\CNF}{\mathsf{CNF}}
\newcommand{\sha}{\operatorname{SHA256}}
\newcommand{\filehash}[1]{\nolinkurl{#1}}

\title{A Certified Order-Twelve Extension of the
\texorpdfstring{$\gamma$--$\theta$}{gamma--theta} Frontier\\
in One-Guard Eternal Domination}
\author{Alec Kriebel}
\date{July 26, 2026}

\begin{document}

\maketitle

\begin{abstract}
In the one-guard-moves eternal domination game, each attack is made at an
unoccupied vertex and exactly one adjacent guard moves to the attacked
vertex.  The $\gamma$--$\theta$ conjecture asserts that
\(\gamma(G)=\gaminf(G)\) implies \(\gamma(G)=\theta(G)\).
MacGillivray, Mynhardt, and Virgile reported an exhaustive computation
excluding counterexamples through order 11.  We extend that published
frontier by proving that no counterexample has order 12.  The order-12
proof has three separately delimited parts: an independently replayed
three-template certificate excludes common parameter \(3\); an exact
18,381-variable formula and a 228,381,671-byte LRAT refutation exclude
connected graphs with common parameter \(4\); and a simplicial
closed-neighborhood reduction, combined with the McCuaig--Shepherd
domination bound, excludes common parameter \(5\).  A classical half-order
characterization proves that these are the only possible parameters.
Consequently, subject explicitly to the published through-order-11
computation, every counterexample has at least 13 vertices.  The campaign
has not independently reproduced the all-graph coverage at orders 10 and
11.  As a structural byproduct, we give an unrestricted, family-level
generalization of a planar independent-antineighborhood lemma of Taletskii.
This finite result neither proves the universal conjecture nor constructs a
counterexample.
\end{abstract}

\section{Introduction}\label{sec:intro}

Let \(G\) be a finite simple graph.  The domination number \(\gamma(G)\)
is the minimum size of a set whose closed neighborhood is all of
\(\V(G)\), and the clique cover number \(\theta(G)\) is the minimum number
of parts in a partition of \(\V(G)\) into cliques.  In the standard
one-guard eternal domination game, guards initially occupy a dominating
set.  After an attack at an unoccupied vertex, exactly one adjacent guard
moves to the attacked vertex, and the new guard set must again be capable
of defending every future attack.  The least sufficient number of guards
is \(\gaminf(G)\).

Klostermeyer and MacGillivray stated the implication
\[
  \gamma(G)=\gaminf(G)\quad\Longrightarrow\quad
  \gamma(G)=\theta(G),
\tag{1.1}
\]
with a proof that was later found to be incomplete
\cite{KlostermeyerMacGillivray2009,KlostermeyerMynhardt2015}.
Klostermeyer and Mynhardt subsequently formulated the question as the
existence problem
\[
  \gamma(G)=\gaminf(G)<\theta(G)
\tag{1.2}
\]
\cite{KlostermeyerMynhardt2015,KlostermeyerMynhardt2016}.  We call (1.1)
the \(\gamma\)--\(\theta\) conjecture.

MacGillivray, Mynhardt, and Virgile reported an exhaustive computation
finding no counterexample through order 11
\cite[Observation 5.6]{MacGillivrayMynhardtVirgile2022}.  Their paper also
catalogs all 56 graphs of order at most 11 for which
\(\gaminf<\theta\).  The present work independently reproduced the
parameters of that catalog and independently enumerated connected graphs
through order 9, but it did not repeat the original cluster-scale
all-graph enumeration at orders 10 and 11.  We therefore keep the
published lower-order computation visible as an explicit premise.

\begin{theorem}[order-12 frontier]\label{thm:frontier}
Assume the published exhaustive result of MacGillivray, Mynhardt, and
Virgile that no counterexample to \emph{(1.1)} has order at most \(11\).
Then no finite simple graph \(G\) of order at most \(12\) satisfies
\[
  \gamma(G)=\gaminf(G)<\theta(G).
\]
Equivalently, every counterexample, if one exists, has order at least \(13\).
\end{theorem}

The order-12 extension in Theorem~\ref{thm:frontier} is assembled from the
following results.

\begin{enumerate}
\item A complete, certificate-backed exclusion of all 12-vertex graphs with
      \(\gamma=\gaminf=3<\theta\) (Theorem~\ref{thm:k3}).
\item A complete, certificate-backed exclusion of connected 12-vertex
      graphs with \(\gamma=\gaminf=4<\theta\)
      (Theorem~\ref{thm:k4}).
\item A structural theorem for simplicial vertices
      (Theorem~\ref{thm:simplicial}) and the McCuaig--Shepherd bound, which
      exclude the order-12 parameter-five case relative to the same
      published lower-order premise (Corollary~\ref{cor:k5}).
\item A classical half-order characterization, which shows that a connected
      order-12 counterexample has common parameter only \(3,4,\) or \(5\)
      (Corollary~\ref{cor:parameters}).
\item An independent-antineighborhood projection theorem for arbitrary
      equality graphs and arbitrary eternal families
      (Theorem~\ref{thm:antineighborhood}), recorded as a generalization of
      Taletskii's planar minimum-counterexample lemma rather than as a
      wholly new antineighborhood idea.
\end{enumerate}

Every computational implication is separated into a graph-to-formula
coverage theorem, an exact byte binding, a proof-certificate replay, and an
independent audit.  Section~\ref{sec:basic} fixes the model and elementary
reductions.  Section~\ref{sec:simplicial} proves the structural reduction.
Sections~\ref{sec:k3} and \ref{sec:k4} give the two certified slices.
Section~\ref{sec:assemble} assembles the frontier, and
Section~\ref{sec:repro} gives reproduction details.

\begin{remark}[scope]\label{rem:scope-intro}
Theorem~\ref{thm:frontier} is a finite extension of a published
computational frontier.  It is not a campaign-only exhaustive enumeration
through order 12, it says nothing about orders at least 13, and it does not
resolve the universal conjecture.  We also make no literature-novelty claim
for the simplicial reduction; one potentially relevant 2018 manuscript
could not be obtained during the audit.  The antineighborhood theorem below
has direct accessible overlap with Taletskii's Lemma 13
\cite{Taletskii2024}, and its prior-art boundary is stated explicitly.
\end{remark}

\section{Model and elementary reductions}\label{sec:basic}

All graphs in this paper are finite, simple, and undirected.  For
\(v\in\V(G)\), write \(N(v)\) and \(N[v]\) for its open and closed
neighborhoods.  A set \(D\subseteq\V(G)\) is dominating if
\(N[D]=\V(G)\).  Let \(\alpha(G)\) denote the independence number and
\(i(G)\) the minimum size of a maximal independent set.  We use
\[
  \theta(G)=\chi(\overline G).
\tag{2.1}
\]

\begin{definition}\label{def:eternal}
An \emph{eternal dominating family of size \(k\)} is a nonempty family
\(\mathcal D\subseteq\binom{\V(G)}k\) such that:
\begin{enumerate}
\item every \(D\in\mathcal D\) dominates \(G\); and
\item for every \(D\in\mathcal D\) and every unoccupied attacked vertex
      \(r\in\V(G)-D\), there is a guard \(u\in D\cap N(r)\) for which
      \[
        (D-\{u\})\cup\{r\}\in\mathcal D.
      \tag{2.2}
      \]
\end{enumerate}
The minimum such \(k\) is \(\gaminf(G)\).
\end{definition}

Thus attacks at occupied vertices are not required, and a response in
(2.2) moves exactly one guard along exactly one edge.  No all-guards-move
parameter is used anywhere in this paper.

\begin{lemma}[independent-set forcing]\label{lem:forcing}
Let \(\mathcal D\) be an eternal dominating family of \(k\)-sets, and let
\(I\) be an independent set.  Starting from any state in \(\mathcal D\),
successive attacks at unoccupied vertices of \(I\) force a state \(D'\) with
\[
  |D'\cap I|=\min\{|I|,k\}.
\]
In particular, every independent \(k\)-set belongs to every eternal
\(k\)-family.
\end{lemma}

\begin{proof}
When an unoccupied \(r\in I\) is attacked, no guard already on \(I\) can
respond because \(I\) is independent.  Every legal response therefore
increases the number of guards on \(I\) by one.  Repeat until no further
such attack is possible.
\end{proof}

\begin{proposition}[parameter chain]\label{prop:chain}
For every nonempty graph \(G\),
\[
  \gamma(G)\leq i(G)\leq\alpha(G)\leq\gaminf(G)\leq\theta(G).
\tag{2.3}
\]
Consequently, if \(\gamma(G)=\gaminf(G)=k\), then
\[
  \gamma(G)=i(G)=\alpha(G)=\gaminf(G)=k,
\tag{2.4}
\]
and \(G\) is well-covered.
\end{proposition}

\begin{proof}
Every maximal independent set dominates, and every maximum independent set
is maximal.  If an eternal family used fewer than \(\alpha(G)\) guards,
Lemma~\ref{lem:forcing} would put all guards on a maximum independent set
while leaving another vertex of it unoccupied; an attack there would have
no adjacent guard.  Thus \(\alpha\leq\gaminf\).

Given a clique partition, place one guard in each part.  An attack is
answered by moving the unique guard within the attacked vertex's clique.
This proves \(\gaminf\leq\theta\).  The equality collapse and
well-coveredness follow by squeezing.
\end{proof}

\begin{lemma}[induced-subgraph monotonicity]\label{lem:induced}
If \(J\) is a nonempty induced subgraph of \(G\), then
\(\gaminf(J)\leq\gaminf(G)\).
\end{lemma}

\begin{proof}
Fix an eternal family in \(G\), let \(m\) be the maximum number of guards
that any family state places in \(\V(J)\), and retain the states attaining
\(m\).  An attack inside \(J\) cannot be answered by moving a guard from
outside \(J\), since the successor would contain more than \(m\) guards in
\(J\).  Thus every response stays within \(J\), preserves \(m\), and
projects to a legal response in the induced graph \(J\).  The projected
states dominate \(J\): each unoccupied vertex has an adjacent responding
guard in \(J\), while occupied vertices dominate themselves.  Also
\(m>0\), since otherwise any attack in the nonempty graph \(J\) would force
a guard into \(J\).  The projections therefore form an eternal family in
\(J\).
\end{proof}

\begin{lemma}[component additivity]\label{lem:additivity}
If \(G_1,\ldots,G_t\) are the nonempty components of \(G\), then
\[
 \gamma(G)=\sum_j\gamma(G_j),\qquad
 \gaminf(G)=\sum_j\gaminf(G_j),\qquad
 \theta(G)=\sum_j\theta(G_j).
\tag{2.5}
\]
\end{lemma}

\begin{proof}
The assertions for \(\gamma\) and \(\theta\) follow by restriction and
union.  Products of componentwise eternal families prove the upper bound
for \(\gaminf\).

For the reverse bound, fix a guard-count vector
\((|D\cap\V(G_1)|,\ldots,|D\cap\V(G_t)|)\) occurring in an eternal family
and retain all states with that vector.  A one-edge move cannot cross
components, so this subfamily remains closed.  Projection to each component
is an eternal family with the corresponding number of guards.  Summing the
componentwise lower bounds proves equality.
\end{proof}

\begin{corollary}[connected reduction]\label{cor:connected}
If a counterexample to \emph{(1.1)} exists, then a connected counterexample
exists.  Every minimum-order counterexample is connected.
\end{corollary}

\begin{proof}
If \(\gamma(G)=\gaminf(G)\), additivity and
\(\gamma(G_j)\leq\gaminf(G_j)\) force equality on every component.  If also
\(\gaminf(G)<\theta(G)\), at least one component has the same strict gap.
\end{proof}

We next record why the common parameter of a counterexample is at least
three.  The following two small values use exactly the model of
Definition~\ref{def:eternal}.

\begin{lemma}\label{lem:odd-antihole}
For every odd \(\ell\geq5\),
\[
  \gaminf(\overline{C_\ell})=3.
\tag{2.6}
\]
\end{lemma}

\begin{proof}
Since \(\theta(\overline{C_\ell})=\chi(C_\ell)=3\), the upper bound follows
from (2.3).  Label \(C_\ell\) cyclically by \(0,\ldots,\ell-1\).
A pair fails to dominate \(\overline{C_\ell}\) exactly when its vertices
have cyclic distance two.

Suppose an eternal two-family existed.  The independent pair \(\{0,1\}\)
belongs to it by Lemma~\ref{lem:forcing}.  From \(\{0,d\}\), for odd
\(1\leq d\leq\ell-6\), attack \(d+2\).  Moving the guard at \(0\) leaves
the nondominating pair \(\{d,d+2\}\), so closure forces
\(\{0,d+2\}\).  Iteration reaches \(\{0,\ell-4\}\).  Attacking
\(\ell-2\) then leaves a cyclic-distance-two pair whichever guard moves.
For \(\ell=5\), this last attack applies to the initial pair.
\end{proof}

\begin{lemma}\label{lem:c7}
\(\gaminf(C_7)=4\).
\end{lemma}

\begin{proof}
A partition into three edges and one singleton gives the upper bound four.
Suppose an eternal family of dominating triples existed.  A dominating
triple has cyclic gap multiset \(\{2,1,1\}\) or \(\{2,2,0\}\).  A triple of
the second type is equivalent to \(\{0,1,4\}\); attacking \(3\) forces the
guard at \(4\) to move, after which vertex \(5\) is undominated.  Hence no
second-type triple belongs to the family.

The maximum independent set \(\{0,2,4\}\) belongs to the family by
Lemma~\ref{lem:forcing}.  Attack \(1\).  Moving the guard at \(0\) leaves
vertex \(6\) undominated, while moving the guard at \(2\) produces the
forbidden triple \(\{0,1,4\}\), a contradiction.
\end{proof}

\begin{proposition}[minimum parameter]\label{prop:min-k}
Every counterexample satisfies
\[
  \gamma(G)=\alpha(G)=\gaminf(G)\geq3.
\tag{2.7}
\]
\end{proposition}

\begin{proof}
The equalities follow from Proposition~\ref{prop:chain}.  Common value one
forces \(\alpha=1\), so \(G\) is complete and \(\theta=1\).

Suppose the common value is two and put \(H=\overline G\).  Then
\(\omega(H)=2<\chi(H)\).  By the Strong Perfect Graph Theorem
\cite{ChudnovskyRobertsonSeymourThomas2006}, \(H\) contains an induced odd
hole or odd antihole.  An odd hole in \(H\) induces an odd antihole in
\(G\), contradicting Lemmas~\ref{lem:induced} and
\ref{lem:odd-antihole}.  An odd antihole in \(H\) induces an odd hole in
\(G\).  At length at least seven this contradicts \(\alpha(G)=2\); at
length five it again contradicts Lemma~\ref{lem:odd-antihole}, since
\(C_5\) is self-complementary.
\end{proof}

\subsection{The half-order restriction}

We use the characterization, due independently to Payan--Xuong and to
Fink--Jacobson--Kinch--Roberts, of graphs without isolated vertices whose
domination number is half their order
\cite{PayanXuong1982,FinkJacobsonKinchRoberts1985}: every component is
either \(C_4\) or a corona \(F\circ K_1\), where one leaf is attached to
each vertex of a connected graph \(F\).

\begin{proposition}\label{prop:half-order}
If \(G\) has no isolated vertices and
\(\gamma(G)=|\V(G)|/2\), then \(\gamma(G)=\theta(G)\).
\end{proposition}

\begin{proof}
For \(C_4\), both parameters are two.  If \(Q=F\circ K_1\) and
\(|\V(F)|=m\), every dominating set must contain a member of each
support--leaf pair, while all supports form a dominating \(m\)-set.
Hence \(\gamma(Q)=m\).  The \(m\) support--leaf edges partition \(Q\) into
cliques, so \(\theta(Q)\leq m\), while the \(m\) leaves are independent and
give \(\theta(Q)\geq m\).  Thus the two parameters agree on every component
in the characterization, and additivity completes the proof.
\end{proof}

\begin{corollary}\label{cor:half-order}
If a connected graph \(G\) of order \(n\geq2\) is a counterexample with
common parameter \(k\), then
\[
  n\geq2k+1.
\tag{2.8}
\]
\end{corollary}

\begin{proof}
A spanning tree of \(G\) is bipartite, and its smaller color class is a
dominating set of the tree and hence of \(G\).  Thus
\(\gamma(G)\leq n/2\).  Equality would contradict
Proposition~\ref{prop:half-order}, so \(k<n/2\).  Integrality gives (2.8).
\end{proof}

\begin{corollary}\label{cor:parameters}
A connected order-12 counterexample has common parameter
\[
  k\in\{3,4,5\}.
\tag{2.9}
\]
\end{corollary}

\begin{proof}
Proposition~\ref{prop:min-k} gives \(k\geq3\), and
Corollary~\ref{cor:half-order} gives \(12\geq2k+1\), hence \(k\leq5\).
\end{proof}

\section{The simplicial closed-neighborhood reduction}
\label{sec:simplicial}

A vertex \(v\) is \emph{simplicial} if \(N[v]\) is a clique.  The next
theorem strengthens the leaf reduction needed later.  Its proof is included
both for completeness and to expose where the one-guard dynamics enter.

\begin{theorem}[simplicial reduction]\label{thm:simplicial}
Let \(G\) satisfy
\[
  \gamma(G)=\gaminf(G)=k,
\]
let \(v\) be simplicial, and suppose \(Q=G-N[v]\) is nonempty.  Then \(Q\)
is well-covered and
\[
  \gamma(Q)=\alpha(Q)=\gaminf(Q)=k-1,
\tag{3.1}
\]
while
\[
  \theta(G)=\theta(Q)+1.
\tag{3.2}
\]
\end{theorem}

\begin{proof}
By Proposition~\ref{prop:chain}, \(G\) is well-covered with
\(\gamma(G)=\alpha(G)=k\).  If \(I\) is a maximal independent set of \(Q\),
then \(I\cup\{v\}\) is a maximal independent set of \(G\): the vertex \(v\)
blocks \(N(v)\), and \(I\) blocks \(Q-I\).  Hence every maximal independent
set of \(Q\) has size \(k-1\), and
\[
  i(Q)=\alpha(Q)=k-1.
\tag{3.3}
\]
In particular, \(\gamma(Q)\leq k-1\).  If a set \(A\) of size at most
\(k-2\) dominated \(Q\), then \(A\cup\{v\}\) would dominate \(G\),
contradicting \(\gamma(G)=k\).  Therefore
\(\gamma(Q)=k-1\).

Fix an eternal \(k\)-family \(\mathcal D\) in \(G\), and choose a maximum
independent set \(I\) of \(Q\).  Lemma~\ref{lem:forcing} places
\(I\cup\{v\}\) in \(\mathcal D\).  Define
\[
  \mathcal E=\{D-\{v\}:D\in\mathcal D,\ v\in D\}.
\tag{3.4}
\]
This family is nonempty.  No state \(D\in\mathcal D\) containing \(v\) can
also contain a vertex \(u\in N(v)\).  If it did, \(D-\{v\}\) would still
dominate \(G\): because \(N[v]\) is a clique, \(u\) dominates everything
that \(v\) could dominate, while \(v\) has no neighbor in \(Q\).  This
would contradict \(\gamma(G)=k\).

Thus each member of \(\mathcal E\) is a \((k-1)\)-subset of \(Q\), and it
dominates \(Q\).  Let \(B=D-\{v\}\in\mathcal E\) and attack
\(r\in\V(Q)-B\).  In \(D=B\cup\{v\}\), the responding guard cannot be
\(v\), since \(v\) has no neighbor in \(Q\).  The response therefore keeps
\(v\) occupied and projects to
\[
  (B-\{u\})\cup\{r\}\in\mathcal E.
\]
Consequently \(\mathcal E\) is an eternal \((k-1)\)-family in \(Q\).
Together with (3.3) and \(\alpha\leq\gaminf\), this proves (3.1).

The clique \(N[v]\), together with a minimum clique partition of \(Q\),
gives \(\theta(G)\leq\theta(Q)+1\).  Conversely, start with a minimum
clique partition of \(G\).  Replace all parts meeting \(N[v]\) by the
single clique \(N[v]\) and the nonempty remainders after deleting \(N[v]\).
The old part containing \(v\) has empty remainder, so the replacement does
not increase the number of parts.  By minimality it is again a minimum
partition.  Deleting its \(N[v]\) part leaves a partition of \(Q\) into
\(\theta(G)-1\) cliques.  Hence
\(\theta(Q)\leq\theta(G)-1\), proving (3.2).
\end{proof}

\begin{corollary}[counterexample preservation]\label{cor:simp-preserve}
Under the hypotheses of Theorem~\ref{thm:simplicial}, if
\[
  \gamma(G)=\gaminf(G)=k<\theta(G),
\]
then \(Q=G-N[v]\) is a smaller counterexample with common parameter \(k-1\).
\end{corollary}

\begin{proof}
Theorem~\ref{thm:simplicial} gives
\(\gamma(Q)=\gaminf(Q)=k-1\) and
\(\theta(Q)=\theta(G)-1\geq k\).
\end{proof}

\begin{corollary}[leaf--support reduction]\label{cor:leaf}
Let \(x\) be a leaf with unique neighbor \(y\), and suppose
\(G-\{x,y\}\) is nonempty.  If
\(\gamma(G)=\gaminf(G)=k\), then \(G-\{x,y\}\) is well-covered,
\[
 \gamma(G-\{x,y\})=\alpha(G-\{x,y\})
 =\gaminf(G-\{x,y\})=k-1,
\]
and
\[
 \theta(G)=\theta(G-\{x,y\})+1.
\]
A strict \(\gamma\)--\(\theta\) gap is preserved.
\end{corollary}

\begin{proof}
The leaf \(x\) is simplicial and \(N[x]=\{x,y\}\), so this is
Theorem~\ref{thm:simplicial} and
Corollary~\ref{cor:simp-preserve}.
\end{proof}

\begin{corollary}\label{cor:no-simplicial}
Every minimum-order counterexample is connected, has no simplicial vertex,
and has minimum degree at least two.  Moreover, the two neighbors of every
degree-two vertex are nonadjacent.
\end{corollary}

\begin{proof}
Connectedness is Corollary~\ref{cor:connected}.  A simplicial vertex would
give a smaller counterexample by Corollary~\ref{cor:simp-preserve}, unless
\(G-N[v]\) were empty; in that case \(G=N[v]\) would be complete.  Hence
there is no simplicial vertex.  In a connected nontrivial graph, a leaf is
simplicial, so the minimum degree is at least two.  A degree-two vertex
whose neighbors are adjacent is also simplicial.
\end{proof}

McCuaig and Shepherd proved that every connected \(n\)-vertex graph \(R\)
with \(\delta(R)\geq2\), except for seven graphs of orders four and seven,
satisfies
\[
  \gamma(R)\leq\frac{2n}{5}
\tag{3.5}
\]
\cite{McCuaigShepherd1989}.  The exception inventory is explicitly
restated as \(\mathcal F_4\cup\mathcal F_7\) by Henning, Schiermeyer, and
Yeo \cite[Theorem 1]{HenningSchiermeyerYeo2011}.

\begin{theorem}[minimum-order bound]\label{thm:five-halves}
Assume the published absence of counterexamples through order \(11\).
If \(G\) is a minimum-order counterexample of order \(n\), with
\[
  k=\gamma(G)=\gaminf(G),
\]
then
\[
  n\geq\left\lceil\frac{5k}{2}\right\rceil.
\tag{3.6}
\]
\end{theorem}

\begin{proof}
The published premise gives \(n\geq12\), so \(G\) is not one of the seven
exceptions of orders four and seven.  Corollary~\ref{cor:no-simplicial}
gives connectedness and minimum degree at least two.  Apply (3.5) and
rearrange:
\[
  k=\gamma(G)\leq\frac{2n}{5}.
\]
Integrality gives (3.6).
\end{proof}

\begin{corollary}[parameter five at order 12]\label{cor:k5}
Assume the published absence of counterexamples through order \(11\).
There is no order-12 graph satisfying
\[
  \gamma(G)=\gaminf(G)=5<\theta(G).
\]
\end{corollary}

\begin{proof}
Such a graph would have minimum possible order, so
Theorem~\ref{thm:five-halves} would give
\[
  12\geq\left\lceil\frac{25}{2}\right\rceil=13,
\]
a contradiction.
\end{proof}

\subsection{A general independent-antineighborhood projection}

The simplicial hypothesis is needed for the clique-cover identity (3.2),
but not for projecting the three equality parameters or the one-guard
strategy.  For \(A\subseteq\V(G)\), write
\[
  N[A]=\bigcup_{a\in A}N[a].
\]

\begin{theorem}[independent-antineighborhood projection]
\label{thm:antineighborhood}
Let \(\gamma(G)=\gaminf(G)=k\), let \(A\) be an independent set of size
\(t<k\), and put \(Q=G-N[A]\).  Then \(Q\) is nonempty and well-covered,
and
\[
  \gamma(Q)=\alpha(Q)=\gaminf(Q)=k-t.
\tag{3.7}
\]
Moreover, every eternal \(k\)-family on \(G\) has an explicit restricted
slice that is an eternal \((k-t)\)-family on \(Q\).
\end{theorem}

\begin{proof}
By Proposition~\ref{prop:chain}, every maximal independent set of \(G\)
has size \(k\).  If \(I\) is any maximal independent set of \(Q\), then
\(I\cup A\) is maximal in \(G\): \(A\) blocks \(N[A]-A\), and \(I\)
blocks \(Q-I\).  Hence every maximal independent set of \(Q\) has size
\(k-t\).  Thus \(Q\) is nonempty and well-covered, with
\[
  i(Q)=\alpha(Q)=k-t.
\tag{3.8}
\]
If \(B\) dominated \(Q\) with at most \(k-t-1\) vertices, then
\(A\cup B\) would dominate \(G\) with at most \(k-1\) vertices.
Consequently \(\gamma(Q)=k-t\).

Fix an arbitrary eternal \(k\)-family \(\mathcal D\) on \(G\).  Choose a
maximum independent set \(I\) of \(Q\).  Lemma~\ref{lem:forcing} gives
\(A\cup I\in\mathcal D\), so
\[
 \mathcal E=
 \{D-A:D\in\mathcal D,\ A\subseteq D,\ D-A\subseteq\V(Q)\}
\tag{3.9}
\]
is nonempty.  If \(C\in\mathcal E\), then \(A\cup C\) dominates \(G\).
No member of \(A\) has a neighbor in \(Q\), so \(C\) dominates \(Q\).

For an arbitrary unoccupied attack \(r\in\V(Q)-C\), closure of
\(\mathcal D\) supplies a guard \(u\in(A\cup C)\cap N(r)\) whose one-edge
move to \(r\) stays in \(\mathcal D\).  Again no member of \(A\) is
adjacent to \(r\), so \(u\in C\); the successor retains \(A\) and all
other guards remain in \(Q\).  It therefore projects to
\[
  (C-\{u\})\cup\{r\}\in\mathcal E.
\]
This proves the exact
\(\forall C\,\forall r\,\exists u\) closure condition on \(Q\), using one
guard and one edge.  Hence \(\gaminf(Q)\leq k-t\), while
\(\alpha(Q)\leq\gaminf(Q)\) and (3.8) give equality.
\end{proof}

\begin{corollary}\label{cor:antineighborhood-minimal}
If \(G\) is a minimum-order counterexample with common parameter \(k\), then
for every nonempty independent \(t\)-set \(A\), \(t<k\),
\[
 \gamma(G-N[A])=\alpha(G-N[A])=\gaminf(G-N[A])
 =\theta(G-N[A])=k-t.
\tag{3.10}
\]
Equivalently, if \(H=\overline G\), then for every \(t\)-clique \(A\) of
\(H\),
\[
 \chi\!\left(H[N_H(A)]\right)
 =\omega\!\left(H[N_H(A)]\right)=k-t,
\tag{3.11}
\]
where \(N_H(A)=\bigcap_{a\in A}N_H(a)\).
\end{corollary}

\begin{proof}
Theorem~\ref{thm:antineighborhood} gives every equality in (3.10) except
the one involving \(\theta\).  If
\(\theta(G-N[A])>k-t\), the proper induced graph \(G-N[A]\) would be a
smaller counterexample.  Minimality gives equality.  Complementation
translates \(G-N[A]\) to \(H[N_H(A)]\), proving (3.11).
\end{proof}

\begin{remark}[prior-art and strength boundary]\label{rem:projection-prior}
Taletskii's Lemma 13 already proves essentially (3.10) for a minimum-order
planar counterexample \cite{Taletskii2024}.  Theorem
\ref{thm:antineighborhood} removes planarity and
the hypothesis that the graph is a minimum counterexample from the
\(\gamma,\alpha,\gaminf\) projection and explicitly projects every eternal
family.  It is therefore presented as an unrestricted, family-level
generalization, not as a wholly new antineighborhood idea.  Its full
independent set form is also obtained by iterating the \(t=1\) case.
For \(k=3\), (3.11) is exactly local bipartiteness together with the
well-covered clique structure, so it gives no new pruning beyond the
odd-wheel obstruction already used in the parameter three analysis.
Novelty relative to the unavailable 2018 manuscript remains unresolved.
\end{remark}

\section{The certified order-12 parameter-three slice}\label{sec:k3}

\begin{theorem}\label{thm:k3}
There is no finite simple graph \(G\) of order \(12\) satisfying
\[
  \gamma(G)=\gaminf(G)=3<\theta(G).
\tag{4.1}
\]
This assertion includes disconnected graphs.
\end{theorem}

We summarize the complete coverage proof and exact certificates.  Full
proof details and all bindings are in the accompanying archive.

\subsection{The exhaustive complement split}

Call a vertex outside an induced cycle a \emph{hub} if it is adjacent to
every rim vertex.  A cycle is \emph{hub-free} if it has no hub.

\begin{proposition}[odd-wheel obstruction]\label{prop:wheel}
If \(\gaminf(G)=3\), then \(\overline G\) contains no induced odd wheel.
\end{proposition}

\begin{proof}
If \(\overline G\) induces \(K_1\vee C_{2q+1}\), then \(G\) induces
\(K_1\mathbin{\dot\cup}\overline{C_{2q+1}}\).
Lemmas~\ref{lem:additivity} and \ref{lem:odd-antihole} give eternal
domination number four for this induced subgraph, contradicting
Lemma~\ref{lem:induced}.
\end{proof}

\begin{proposition}[three-template split]\label{prop:k3-split}
If \(G\) satisfies (4.1), then \(G\) is connected and
\(H=\overline G\) contains a hub-free induced \(C_\ell\) for some
\[
  \ell\in\{5,7,9\}.
\tag{4.2}
\]
\end{proposition}

\begin{proof}
If \(G\) were disconnected, component additivity would give a
counterexample component.  Proposition~\ref{prop:min-k} says that component
has parameter at least three, so it consumes the entire domination budget;
hence \(G\) is connected.

Proposition~\ref{prop:chain} gives
\[
  \omega(H)=\alpha(G)=3<\chi(H)=\theta(G).
\tag{4.3}
\]
The Strong Perfect Graph Theorem supplies an induced odd hole or odd
antihole.  An odd antihole of length \(2q+1\) has clique number \(q\), so
its length is five or seven.  The five-antihole is \(C_5\).  An induced
\(\overline{C_7}\) in \(H\) induces \(C_7\) in \(G\), contradicting
Lemmas~\ref{lem:induced} and \ref{lem:c7}.  Thus \(H\) contains an odd hole.

It is hub-free by Proposition~\ref{prop:wheel}.  Since no pair dominates
\(G\), every pair in \(H\) has a common \(H\)-neighbor.  The endpoints of
a rim edge have no common neighbor on an induced rim, so they have one
outside it.  If only one vertex lay outside the hole, it would be a hub.
At least two vertices therefore lie outside, and the odd rim has length at
most ten.  This gives (4.2).
\end{proof}

\subsection{Exact formulas and coverage}

Fix \(\ell\in\{5,7,9\}\), label an induced cycle of
\(H=\overline G\) by \(0,\ldots,\ell-1\), and label one common
\(H\)-neighbor of rim edge \(01\) as vertex \(\ell\).  The formula has
6,886 variables:
\[
\begin{array}{rcll}
66   &:& e_{uv} & (uv\in E(H)),\\
660  &:& w_{ab,c} & (c\text{ witnesses a common }H\text{-neighbor of }a,b),\\
220  &:& f_D & (D\in\binom{[12]}3\text{ is selected}),\\
5,940&:& m_{D,r,u} & (u\in D\text{ responds to }r\notin D).
\end{array}
\tag{4.4}
\]
The clauses assert: no \(K_4\) in \(H\); a common \(H\)-neighbor for every
pair; the labeled hub-free cycle template; connectedness of \(G\); a
nonempty selected family of dominating triples; one legal successor for
every selected state and every unoccupied attack; and selection of every
\(H\)-triangle.  In particular, a move literal implies
\(\neg e_{ur}\), because a legal move uses an edge of \(G\), not \(H\).
It also selects the state \((D-\{u\})\cup\{r\}\), whose own clauses require
domination.

For every normalized map \(c:[12]\to\{0,1,2\}\) proper on the forced
template, the coloring obstruction is
\[
  C_c=\bigvee_{\substack{u<v\\c(u)=c(v)}}e_{uv}.
\tag{4.5}
\]
It is false exactly when \(c\) is a proper coloring of \(H\).
The \(C_5\) and \(C_7\) formulas contain complete banks of respectively
3,645 and 1,701 such clauses.  The \(C_9\) formula contains 170 valid
coloring clauses obtained during counterexample-guided refinement.  No
completeness claim is needed for that subset: every graph with
\(\chi(H)>3\) satisfies every clause (4.5).

The \(C_5\) branch also lexicographically orders the six undistinguished
outer signatures.  The formula is invariant under their full permutation
group, so every model has an orbit representative satisfying the order.

\begin{lemma}[parameter-three graph-to-formula coverage]
\label{lem:k3-coverage}
Every graph satisfying (4.1) induces a satisfying assignment of at least
one of the three exact formulas in Table~\ref{tab:k3}.
\end{lemma}

\begin{proof}
Use Proposition~\ref{prop:k3-split} and label the resulting template as
above.  Set edge and witness variables from \(H\).
Equation (4.3) supplies the no-\(K_4\) and coloring clauses; failure of every
pair to dominate supplies the common-neighbor clauses; and connectedness
and hub-freeness supply the corresponding template clauses.

Set \(f_D\) on an actual nonempty eternal three-family and choose one
one-guard response for each selected state and unoccupied attack.
Lemma~\ref{lem:forcing} selects every \(H\)-triangle.  If \(\ell=5\),
permute undistinguished vertices and all associated semantic variables to
satisfy the signature order.  Thus the exact branch formula is satisfied.
\end{proof}

\subsection{Certificates}

\begin{table}[ht]
\centering
\small
\caption{Exact parameter-three branch formulas and accepted
addition-only RUP proofs.}
\label{tab:k3}
\begin{tabular}{@{}lrrr@{}}
\toprule
 & \(C_5\) & \(C_7\) & \(C_9\)\\
\midrule
variables & 6,886 & 6,886 & 6,886\\
clauses & 23,968 & 21,718 & 20,200\\
literals & 192,169 & 148,551 & 117,841\\
CNF bytes & 754,323 & 621,864 & 530,053\\
proof additions & 247,981 & 284,317 & 4,705\\
proof bytes & 6,337,621 & 18,093,724 & 65,906\\
\bottomrule
\end{tabular}
\end{table}

In \(C_5,C_7,C_9\) order, the formula SHA-256 values are
\begin{center}
\begin{minipage}{0.95\linewidth}
\filehash{c6a0811c718ff8e9352f253e4ce225ce2826def9c3e4a9cd55f0b0152703d104}\\
\filehash{6a011e685e58ef517f2ab8253ca40987bd7b742a470bedbacdc3a5e94fc995a7}\\
\filehash{2845f242a094484a8d114e70ca1a8678dfcff79fadd56bd57813e25c2e49523d}
\end{minipage}
\end{center}
The corresponding addition-only proof SHA-256 values are
\begin{center}
\begin{minipage}{0.95\linewidth}
\filehash{c6c24853e30073e66fb396441edb176a0160d062a8558e25fa18a955f33927c3}\\
\filehash{e8052df40d3e0c39b945a8735889039daba55eacc351e1822828b3d94f7baae9}\\
\filehash{24c5647d3a57f2de221fba96747c618575a3aba086c5e4bca17aade55ce7d4ab}
\end{minipage}
\end{center}

Each proof is an addition-only sequence of reverse-unit-propagation (RUP)
consequences ending in the empty clause, so it is a directly checkable
unsatisfiability certificate \cite{WetzlerHeuleHunt2014}.  Strict replays
reported zero RAT lemmas.  The \(C_5\) binary proof was independently
parsed and replayed; the \(C_7\) proof passed two warning-fatal RUP-only
replays; and the \(C_9\) proof passed both a separately written
standard-library checker and a pinned external checker.  Two independent
reviews audited the three-branch coverage theorem.

\begin{proof}[Proof of Theorem~\ref{thm:k3}]
A putative graph gives a satisfying assignment of one exact formula by
Lemma~\ref{lem:k3-coverage}.  The corresponding accepted RUP refutation
proves that formula unsatisfiable, a contradiction.
\end{proof}

\section{The certified connected order-12 parameter-four slice}
\label{sec:k4}

\begin{theorem}\label{thm:k4}
There is no connected finite simple graph \(G\) of order \(12\) satisfying
\[
  \gamma(G)=\gaminf(G)=4<\theta(G).
\tag{5.1}
\]
\end{theorem}

Put \(H=\overline G\), and let \(e_{uv}=1\) mean \(uv\in E(H)\).
Equality collapse gives
\[
  \gamma(G)=\alpha(G)=\gaminf(G)=4<\theta(G).
\tag{5.2}
\]

\subsection{An exact anchored synthesis formula}

\begin{proposition}[static characterization]\label{prop:k4-static}
Up to relabeling vertices, \(\gamma(G)=\alpha(G)=4\) if and only if:
\begin{enumerate}
\item \(A=\{0,1,2,3\}\) is a \(K_4\) in \(H\);
\item \(H\) has no \(K_5\); and
\item every three-set \(S\) has a common \(H\)-neighbor outside \(S\).
\end{enumerate}
\end{proposition}

\begin{proof}
If \(\gamma=\alpha=4\), label a maximum independent four-set of \(G\) as
\(A\).  It is a \(K_4\) of \(H\), and \(\alpha(G)=4\) forbids a \(K_5\)
in \(H\).  No three-set dominates \(G\), so every three-set has a vertex
adjacent in \(G\) to none of it, equivalently a common \(H\)-neighbor.

Conversely, the anchor and no-\(K_5\) clauses give \(\alpha(G)=4\).
The common-neighbor condition says that no three-set dominates \(G\), so
\(\gamma(G)\geq4\).  The independent four-set \(A\) is maximum and hence
maximal, so it dominates \(G\) and \(\gamma(G)\leq4\).
\end{proof}

The exact formula uses the following 18,381 variables.
\[
\begin{array}{@{}lrr@{}}
\toprule
\text{role} & \text{family} & \text{count}\\
\midrule
\text{complement edges} & e_{uv},\ u<v & 66\\
\text{triple common-neighbor witnesses}
  & w_{S,x},\ |S|=3,\ x\notin S & 1,980\\
\text{selected guard states} & f_D,\ |D|=4 & 495\\
\text{one-guard responses}
  & m_{D,r,u},\ r\notin D,\ u\in D & 15,840\\
\midrule
\text{total}&&18,381\\
\bottomrule
\end{array}
\tag{5.3}
\]
Its static clauses encode Proposition~\ref{prop:k4-static} and
connectedness of \(G\).  For every nonempty proper cut \(S\) containing
vertex \(0\), the connectedness clause is
\[
  \bigvee_{\substack{u\in S\\v\notin S}}\neg e_{uv},
\tag{5.4}
\]
where a negative \(H\)-edge is a \(G\)-edge.

For each four-set \(D\) and \(x\notin D\), domination is
\[
  \neg f_D\ \vee\ \bigvee_{u\in D}\neg e_{ux}.
\tag{5.5}
\]
At least one \(f_D\) is true.  For each selected \(D\) and unoccupied
attack \(r\notin D\), the clauses require at least one \(m_{D,r,u}\), and
each true response implies
\[
  \neg e_{ur}
  \quad\text{and}\quad
  f_{(D-\{u\})\cup\{r\}}.
\tag{5.6}
\]
Thus the encoded quantifier is exactly
\[
  \forall D\in\mathcal D\ \forall r\notin D\ \exists u\in D,
\]
with one named guard moving along one \(G\)-edge to a selected successor.
Every \(H\)-\(K_4\) is also selected, a sound strengthening by
Lemma~\ref{lem:forcing}.

The anchored \(K_4\) fixes four distinct colors in any proper
four-coloring of \(H\).  Normalize them by \(c(j)=j\) for \(j\in A\).
For each of the \(4^8=65,536\) colorings of the remaining vertices, append
\[
  C_c=\bigvee_{\substack{u<v\\c(u)=c(v)}}e_{uv}.
\tag{5.7}
\]
The full bank is equivalent to \(\chi(H)>4\), hence to \(\theta(G)>4\).

Finally, lexicographically order the four-bit anchor signatures
\[
  (e_{0v},e_{1v},e_{2v},e_{3v}),\qquad v=4,\ldots,11.
\tag{5.8}
\]
Permuting the eight outer vertices and all semantic variables preserves
the base formula and permutes the complete coloring bank, so every orbit
has an ordered representative.

\begin{proposition}[parent coverage]\label{prop:k4-parent}
The resulting anchored parent formula is satisfiable if and only if, up to
relabeling one maximum independent four-set, there is a connected
12-vertex graph satisfying (5.2).
\end{proposition}

\begin{proof}
Given such a graph, Proposition~\ref{prop:k4-static} supplies the static
assignment, an actual eternal family supplies (5.5)--(5.6), and
\(\chi(H)>4\) supplies (5.7).  Relabel outer vertices to satisfy (5.8).

Conversely, the static clauses give \(\gamma=\alpha=4\), the selected-state
clauses give an eternal four-family and hence \(\gaminf=4\), the cut
clauses give connectedness, and the complete bank gives \(\theta>4\).
\end{proof}

The exact parent DIMACS file has
\[
\begin{array}{rcl}
\text{variables}&=&18,381,\\
\text{clauses}&=&114,742,\\
\text{literal occurrences}&=&1,180,016,\\
\text{bytes}&=&3,992,947.
\end{array}
\tag{5.9}
\]
Its SHA-256 is
\begin{center}
\filehash{adbe0c01614bae6cd3aed4ccdcd45a757ca56e7ef9c4f2f280f2d8ef200e40ac}.
\end{center}

\subsection{DoubleLex orbit reduction}

Form the \(8\)-by-\(4\) anchor--outer incidence matrix
\[
  M_{v,j}=e_{jv},\qquad v\in\{4,\ldots,11\},\ j\in A.
\tag{5.10}
\]
The group \(S_8\times S_4\) permutes its rows and columns.  Relabeling
vertices simultaneously in all graph, witness, family, and move variables
preserves the formula.  Under an anchor permutation, permuting the four
color names restores the normalized anchor colors, so the complete bank is
also invariant.

\begin{proposition}[DoubleLex equisatisfiability]
\label{prop:doublelex}
The parent formula is satisfiable if and only if it has a model in which
both the rows and columns of \(M\) are lexicographically nondecreasing.
\end{proposition}

\begin{proof}
Choose, among the finite \(S_8\times S_4\) orbit of a parent model, an
image whose row-major matrix word is lexicographically least.  If an
earlier row exceeded a later row, swapping those rows would decrease the
word.  Hence the rows are ordered.  If adjacent columns were inverted,
swap them.  At their first differing row, the row-major word changes from
\(1\) to \(0\), with all earlier entries fixed, again a contradiction.
Thus the same least image has both orders.  The reverse implication is
immediate.
\end{proof}

Three auxiliary-free eight-bit comparators impose the column order.  They
add 765 clauses and 10,758 literals.  The exact strengthened formula \(D\)
therefore has
\[
\begin{array}{rcl}
\text{variables}&=&18,381,\\
\text{clauses}&=&115,507,\\
\text{literal occurrences}&=&1,190,774,\\
\text{bytes}&=&4,030,657.
\end{array}
\tag{5.11}
\]
The exact SHA-256 binding is
\begin{center}
\filehash{14284db1f0b9cfb37b91d834fbabac1d0ca06d36e0d2782683e35cbd04a976e7}.
\end{center}
An independent reconstruction reproduced both the complete parent prefix
and the 37,710-byte comparator suffix byte for byte.  Exhaustive comparator
tests and semantic mutation tests checked the symmetry and one-guard
encoding.

\subsection{The exact LRAT refutation}

A complete solver proof was normalized to an addition-only RUP stream of
15,783,377 bytes with SHA-256
\begin{center}
\filehash{2741335a5ed9af769f0db4bd0c03a70e414d0568681d5b8261a5667ed30b6686}.
\end{center}
Warning-fatal forward and backward replays both verified the stream and
reported zero RAT lemmas.  Conversion produced a 228,381,671-byte LRAT
certificate with SHA-256
\begin{center}
\filehash{0e04eb639a3f7f7d335126d56040abb4ef11e8548770262c316a608390659263}.
\end{center}
The pinned, separately checked \texttt{lrat-check} executable accepted it.

The independent hostile audit reconstructed \(D\) and the normalized proof
byte for byte, reran both RUP directions, generated a fresh byte-identical
LRAT, replayed both retained and fresh LRATs, and rejected six mutated
artifacts.  Its exact verdict was
\begin{center}
\nolinkurl{ACCEPT_EXACT_DOUBLELEX_CNF_UNSAT_ONLY}.
\end{center}
A separate implication audit checked the transfer from exact \(D\)-UNSAT through
Propositions~\ref{prop:k4-parent} and \ref{prop:doublelex}.

\begin{proof}[Proof of Theorem~\ref{thm:k4}]
If a connected graph satisfied (5.1), equality collapse would give (5.2),
and Proposition~\ref{prop:k4-parent} would give a parent model.
Proposition~\ref{prop:doublelex} would then give a model of \(D\).
The accepted LRAT certificate proves \(D\) unsatisfiable, a contradiction.
\end{proof}

\begin{remark}[connected scope]\label{rem:k4-connected}
Theorem~\ref{thm:k4} is a connected exclusion.  It does not, by itself,
exclude every disconnected parameter-four graph of order 12.  In the
proof of Theorem~\ref{thm:frontier}, the published absence of smaller
counterexamples makes a putative order-12 counterexample minimum-order, and
Corollary~\ref{cor:connected} then supplies connectedness.
\end{remark}

\section{Assembly of the order-12 frontier}\label{sec:assemble}

\begin{proof}[Proof of Theorem~\ref{thm:frontier}]
Orders at most 11 are excluded by the explicitly assumed published
computation.  Suppose an order-12 counterexample \(G\) exists.  It then has
minimum possible order and is connected by
Corollary~\ref{cor:connected}.  Put
\[
  k=\gamma(G)=\gaminf(G).
\]
Corollary~\ref{cor:parameters} leaves exactly
\[
  k\in\{3,4,5\}.
\]
Theorem~\ref{thm:k3} excludes \(k=3\);
Theorem~\ref{thm:k4} excludes the connected \(k=4\) case; and
Corollary~\ref{cor:k5} excludes \(k=5\).  No parameter remains.
\end{proof}

\begin{remark}[evidentiary boundary]
The \(k=3\) and connected \(k=4\) order-12 slices are backed by campaign
certificates and independent replays.  The \(k=5\) slice is analytic,
relative to the classical McCuaig--Shepherd theorem and the published
through-order-11 computation.  The all-parameter statement through order
12 additionally imports that published computation at orders at most 11.
The campaign's audit of the published 56-graph appendix is a regression
check, not a replacement for the original all-graph coverage.
\end{remark}

\section{Certificates, reproduction, and provenance}\label{sec:repro}

All paths below are relative to the archive root
\texttt{gamma\_theta\_eternal\_domination/}.

\subsection{Parameter-three replay}

The accepted theorem note and its two independent composition reviews are:
\begin{center}
\begin{minipage}{0.95\linewidth}
\texttt{math/lemmas/order12\_k3\_exclusion.md}\\
\texttt{reviews/order12\_k3\_exclusion\_hostile\_review.md}\\
\texttt{reviews/order12\_k3\_exclusion\_second\_review.md}.
\end{minipage}
\end{center}
The claim-level acceptance record is
\begin{center}
\texttt{results/order12\_k3\_exclusion\_acceptance.json}.
\end{center}
From the archive root, the fail-closed full replay is:
\begin{verbatim}
python3 -I -B repro/c035/replay.py --mode full \
  --output /fresh/path/c035-full-replay.json
\end{verbatim}
It validates exact source, formula, proof, checker, review, and manifest
bindings before running the three proof replays in fresh private
directories.  A successful record says
\texttt{PASS\_FULL\_C035\_REPLAY} and
\texttt{proofs\_freshly\_replayed=true}.  The fast mode checks metadata
only and explicitly returns no mathematical claim.

\subsection{Parameter-four replay}

The compact publication package is
\begin{center}
\texttt{certificates/order12\_k4\_doublelex\_seed0\_lrat\_publication/}.
\end{center}
It contains the 64,288,636-byte compressed LRAT with SHA-256
\begin{center}
\filehash{edc0f6b76bc96b2b26677f399566ae437cc47a6fc0cc921eaff81d77b72a50da}.
\end{center}
Decompression is accepted only if the result has the size and hash in
Section~\ref{sec:k4}.  From the archive root, the one-command replay is:
\begin{verbatim}
python3 \
  certificates/order12_k4_doublelex_seed0_lrat_publication/\
verify_publication.py
\end{verbatim}
The verifier hashes the formula, package metadata, compressed and recovered
proofs, checker, author package, and independent review before invoking
\texttt{lrat-check}.  Its accepted terminal verdict is
\texttt{VERIFIED\_EXACT\_DOUBLELEX\_CNF\_UNSAT\_ONLY}.  The graph transfer
is separately recorded in:
\begin{center}
\begin{minipage}{0.95\linewidth}
\texttt{math/lemmas/order12\_k4\_synthesis\_target.md}\\
\texttt{math/lemmas/order12\_k4\_doublelex.md}\\
\texttt{reviews/order12\_k4\_doublelex\_conditional\_implication\_audit/}\\
\hspace*{1em}\texttt{REVIEW.md}.
\end{minipage}
\end{center}

\subsection{Frontier and structural reviews}

The assembled theorem and the hash-bound independent review are:
\begin{center}
\begin{minipage}{0.95\linewidth}
\texttt{math/lemmas/order12\_frontier.md}\\
\texttt{reviews/order12\_frontier\_second\_review/REVIEW.md}.
\end{minipage}
\end{center}
Their SHA-256 values are respectively
\begin{center}
\filehash{adb27204d33feb47933f2a4b1e381485b2e1b80c22b56a67b18586c4933c2b75}
\end{center}
and
\begin{center}
\filehash{875447d219e5670d66be7d6e4b7ec9f9b3b03bba4d2b169b836706b624a92926}.
\end{center}
The review's exact verdict is
\texttt{ACCEPT\_ORDER12\_FRONTIER\_WITH\_EXPLICIT\_PUBLISHED\_PREMISE}.
The simplicial theorem was separately audited in
\begin{center}
\texttt{reviews/simplicial\_neighborhood\_reduction\_hostile/}.
\end{center}
That review found no mathematical defect and exhaustively tested the
reduction on all 13,598 graphs in the retained small-graph universe through
order 8.  Finite testing is supporting regression evidence; the proof of
Theorem~\ref{thm:simplicial} is the mathematical evidence.

The independent-antineighborhood theorem, its exact prior-art audit, and
its clean-room falsification probe are:
\begin{center}
\begin{minipage}{0.95\linewidth}
\texttt{math/lemmas/independent\_antineighborhood\_projection.md}\\
\texttt{reviews/general\_neighborhood\_projection\_independent/REVIEW.md}\\
\texttt{reviews/general\_neighborhood\_projection\_independent/probe.py}.
\end{minipage}
\end{center}
The probe imports no campaign evaluator and checks all 13,598 unlabeled
graphs through order 8, including 14,421 eligible independent sets and
56,166 one-guard attack obligations.  Its role is falsification and
regression; the proof of Theorem~\ref{thm:antineighborhood} is the
unbounded-order evidence.

\subsection{Research and verification provenance}

This manuscript and the associated research artifacts were produced with
substantial assistance from autonomous AI systems.  No theorem or finite
claim was accepted on model output alone.  Mathematical reductions were
assigned to separate adversarial proof reviews; exact formulas were
independently reconstructed; proof streams were replayed by checkers that
do not import the search implementation's graph-transition core; and
mutation tests were used to detect complement-sign, occupied-attack,
all-guards-move, missing-successor, and malformed-certificate errors.
These procedures reduce, but do not eliminate, the need for human
inspection.  Human authors remain responsible for any submitted version
and for the archived evidence.

\section*{Data availability}

The formulas, proofs, verifiers, manifests, reviews,
literature-source identifiers, and exact hashes cited above are publicly
archived in the tagged release
\href{https://github.com/AlecKriebel/Math/releases/tag/gamma-theta-order12-frontier-v1.0.0}
{\texttt{gamma-theta-order12-frontier-v1.0.0}}.
The corresponding human-readable
\href{https://aleckriebel.github.io/Math/research/gamma-theta-conjecture/}
{project page} records the active campaign status.

\section*{Acknowledgment of limitations}

No external expert was contacted during this independent campaign.  The
universal conjecture remains open in this work.  The strongest conclusion
is the explicitly delimited order-12 frontier theorem, and its lower-order
part relies on the published computation of
MacGillivray--Mynhardt--Virgile.

\bibliographystyle{plain}
\bibliography{references}

\end{document}
